\chapter{Modèles génératifs : contexte et enjeux}
\label{chap:generative}

\section{Le problème de la génération}

On dispose d'un jeu de données $\{x^{(1)}, \dots, x^{(N)}\}$ — par exemple $N$ photographies de
visages — que l'on suppose échantillonnées indépendamment d'une distribution inconnue $p_{\text{data}}$
sur un espace de très grande dimension (une image $64\times 64$ en couleur vit dans
$\mathbb{R}^{64 \times 64 \times 3} = \mathbb{R}^{12288}$).

\begin{definition}[Modèle génératif]
Un modèle génératif est une procédure qui, sans connaître $p_{\text{data}}$ explicitement, permet de
produire de nouveaux échantillons $x \sim p_\theta$ où $p_\theta$ est une distribution apprise à
partir des données, censée approcher $p_{\text{data}}$.
\end{definition}

La difficulté centrale n'est pas de mémoriser les données d'entraînement (un modèle qui recopie une
image du jeu d'entraînement au hasard \og génère \fg\ techniquement une image plausible, mais n'a rien
appris d'utile) : elle est de capturer la \emph{structure} de $p_{\text{data}}$ suffisamment bien pour
produire des échantillons \textbf{nouveaux} et \textbf{plausibles}. Pour des visages, cela veut dire
respecter implicitement des milliers de contraintes non écrites (deux yeux symétriques, une peau dont
la texture est cohérente avec l'éclairage, des proportions anatomiquement crédibles) sans qu'aucune de
ces contraintes n'ait été programmée à la main.

\section{Un bref panorama des approches}

\subsection{Modèles à variable latente explicite : les VAE}

Un \emph{auto-encodeur variationnel} (VAE) apprend conjointement un encodeur $q_\phi(z \mid x)$, qui
compresse une image $x$ en un vecteur latent $z$ de basse dimension, et un décodeur $p_\theta(x \mid z)$,
qui reconstruit une image à partir de $z$. En imposant que $z$ suive approximativement une loi normale
$\mathcal{N}(0, I)$, on peut, après entraînement, générer une nouvelle image en tirant
$z \sim \mathcal{N}(0, I)$ et en appliquant le décodeur.

\begin{intuition}
Le VAE apprend une \og carte \fg\ compacte de l'espace des images : chaque point de cette carte (le
latent $z$) correspond à une image. Générer, c'est piocher un point au hasard sur la carte.
\end{intuition}

En pratique, les VAE produisent des images souvent légèrement floues : la loss (une borne inférieure
de la vraisemblance, l'ELBO) pénalise moins une image \og moyennée \fg\ qu'une hallucination nette
mais incorrecte, ce qui pousse le modèle vers des reconstructions prudentes.

\subsection{Modèles adversariaux : les GAN}

Un \emph{Generative Adversarial Network} (GAN) oppose deux réseaux : un générateur $G$ qui transforme
du bruit en image, et un discriminateur $D$ qui essaie de distinguer les vraies images des images
générées. Les deux réseaux s'entraînent simultanément, en compétition — le générateur cherchant à
tromper le discriminateur, le discriminateur cherchant à ne pas se laisser tromper.

Les GAN produisent typiquement des images très nettes (pas de flou caractéristique du VAE), mais
l'entraînement est notoirement instable (risque d'effondrement de mode, où le générateur cesse de
produire une variété d'images et se contente de reproduire quelques exemples qui trompent bien le
discriminateur) et ne dispose d'aucune fonction de perte directement interprétable pour suivre la
progression.

\subsection{Modèles à transport continu : diffusion et flow matching}

La famille qui nous intéresse dans ce document adopte une stratégie différente : au lieu d'apprendre
directement une transformation \og bruit $\to$ image \fg\ en une seule étape (comme le générateur d'un
GAN), on apprend une transformation \textbf{progressive}, décomposée en de nombreux petits pas, chacun
faisant passer d'un état légèrement plus bruité à un état légèrement moins bruité (ou l'inverse selon
le sens dans lequel on parcourt le processus).

\begin{intuition}
Sculpter un bloc de marbre en une statue en un seul geste est presque impossible. Le faire par des
milliers de petits coups de burin successifs, chacun corrigeant légèrement la forme précédente, est
beaucoup plus faisable — même si aucun coup de burin individuel ne \og sait \fg\ à quoi ressemblera la
statue finale. C'est l'intuition centrale des modèles de diffusion et de flow matching : remplacer un
problème de génération difficile (un grand saut) par une séquence de problèmes de \emph{raffinement
local} beaucoup plus faciles (de petits pas).
\end{intuition}

Les deux chapitres suivants détaillent chacune de ces deux approches : les modèles de diffusion
(chapitre~\ref{chap:ddpm}), historiquement les premiers à obtenir des résultats compétitifs avec les
GAN sur la génération d'images, puis le flow matching (chapitre~\ref{chap:flow}), une formulation plus
récente et plus directe du même principe.

\begin{table}[h]
\centering
\begin{tabular}{@{}lccc@{}}
\toprule
& \textbf{VAE} & \textbf{GAN} & \textbf{Diffusion / Flow Matching} \\
\midrule
Entraînement & stable & instable & stable \\
Netteté des images & moyenne (flou) & élevée & élevée \\
Diversité des échantillons & bonne & risque d'effondrement & bonne \\
Coût de génération & 1 passe & 1 passe & de dizaines à milliers de passes \\
Loss interprétable & oui (ELBO) & non (jeu à somme nulle) & oui (MSE) \\
\bottomrule
\end{tabular}
\caption{Comparaison schématique des grandes familles de modèles génératifs d'images.}
\end{table}

Le prix à payer pour la stabilité d'entraînement et la qualité des modèles à transport continu est le
\textbf{coût de génération} : produire une seule image nécessite d'appeler le réseau de neurones de
nombreuses fois (typiquement de quelques dizaines à un millier de fois), contre une seule passe pour un
GAN ou un VAE. C'est un compromis assumé par toute cette famille de modèles, diffusion comme flow
matching.
